% ==========================================
% BAB I PENDAHULUAN
% ==========================================
\chapter{PENDAHULUAN}
\label{chap:pendahuluan}
% --- Latar Belakang ---
\section{Latar Belakang}
Sejak pertama kali komputer dikenalkan pada tahun 1960, banyak kelompok yang membuat aplikasi untuk mencatat sumber daya. Pada tahun 1970 \textit{Materials Requirements Planning} (MRP) pertama kali dikenalkan, kemudian berkembang menjadi MRP II dengan sistem manajemen sumber daya yang lebih komperhensif. Aplikasi ini kemudian berkembang dari pencatatan dan pengaturan menjadi alat bantu manajemen operasional yang disebut \textit{Enterprise Resource Planning} (ERP). Dengan ERP \textit{enterprise} dapat melakukan manajemen keuangan, sumber daya manusia, hingga manajemen rantai pasokan dalam satu \textit{platform}. Pada tahun 1990an \textit{market} ERP menjadi sangat kompetitif, namun setelah \textit{dotcom bubble} pada 2001 hanya tersisa tiga \textit{vendor} utama (yang awalnya dari enam atau lebih) yaitu SAP, Oracle, dan Infor \cite{goldston2020evolution}.

Pada awalnya, ERP merupakan sebuah aplikasi untuk mengatur sumber daya pada \textit{enterprise} besar yang membutuhkan biaya investasi awal yang sangat besar. Namun setelah meledaknya popularitas \textit{cloud computing} yang memungkinkan alternatif biaya berulang dan (hampir) tidak ada biaya investasi awal, usaha mikro, kecil, dan menengah (UMKM) mulai tertarik \cite{goldston2020evolution}. Keunggulan utama ERP adalah sifat interoperabilitasnya yang sangat baik, dirancang berdasarkan rantai pasok atau \textit{value chain} \cite{boza2015}. Interoperabilitas baik, \textit{platform} satu untuk semua, dan fleksibilitas \textit{deployment} (\textit{cloud} atau \textit{on-premise}) membuat ERP menjadi \textit{go-to} untuk manajemen sumber daya \textit{enterprise}. Hingga saat ini, popularitas ERP semakin berkembang dengan \textit{vendor} lama semakin populer dan munculnya \textit{vendor} baru yang menawarkan untuk UMKM contohnya Odoo dan Mekari.

Pada umumnya, ERP dapat di-\textit{deploy} menggunakan dua model infrastruktur yaitu \textit{cloud} dan \textit{on-premise}, walaupun belakangan ini muncul alternatif ketiga yaitu \textit{hybrid}. Untuk \textit{cloud} sendiri ada tiga jenis layanan utama yaitu \textit{Infrastructure-as-a-Service} (IaaS), \textit{Platform-as-a-Service} (PaaS), dan \textit{Software-as-a-Service} (SaaS). IaaS pemberi layanan menyediakan infrastruktur seperti \textit{server}, \textit{storage}, atau sumber daya komputasi (contohnya VPS), Paas menyediakan \textit{platform} komputasi tanpa berurusan langsung dengan \textit{operating system} (OS), dan SaaS menyediakan akses aplikasi melalui jaringan internet \cite{odunayo2018cloud}. Model \textit{on-premise}, berarti ERP dijalankan pada \textit{server} yang dimiliki sendiri. Model ini memiliki biaya investasi awal yang tinggi namun tidak memiliki biaya berulang seperti \textit{cloud} (pada umumnya) selain listrik \cite{Younus2024}.

Untuk mengatasi masalah biaya kepemilikan atau biaya berulang tinggi, diperlukan infrastruktur yang hemat biaya dan hemat daya. Salah satu kandidat yang menjanjikan adalah perangkat berbasis prosesor \textit{advanced RISC machine} (ARM). Prosesor ARM menerapkan arsitektur \textit{Reduced Instruction Set Computing} (RISC) yang dalam satu dekade terakhir berkembang dengan pesat. RISC atau ARM dikenal dengan arsitekturnya yang dirancang agar lebih efisien dibandingkan kompetitor utamanya yaitu x86 (arsitektur prosessor terkenal untuk komputer dan \textit{server}) \cite{blem2013power, aroca2012towards}. Efisiensi telah menjadi senjata utama ARM yang membuatnya dominan dalam kategori perangkat seperti \textit{smartphone} dan \textit{tablet} hingga \textit{Single-Board Computer} (SBC) seperti Raspberry Pi dan Orange Pi.

Perkembangan teknologi ARM yang berkelanjutan telah meningkatkan kemampuan komputasional SBC secara signifikan, yang awalnya hanya digunakan sebagai \textit{microcontroller} telah berubah menjadi komputer bahkan \textit{server} \cite{Ariza2022}. Pada "Implementasi \textit{Cluster Server} Pada Raspberry Pi Dengan Menggunakan Metode \textit{Load Balancing}" peneliti berhasil menjalankan \textit{web server} menggunakan Raspberry Pi 4 \cite{Putra2019}. Walaupun berhasil menjalankan \textit{web server}, penelitian tersebut hanya menguji skenario \textit{deployment} yang relatif sederhana—spesifiknya \textit{web server} dengan \textit{traffic} HTTP—yang tidak representative terhadap kompleksitas aplikasi ERP. ERP melibatkan operasi \textit{database} yang intensif, manajemen \textit{state} yang kompleks, pemrosesan \textit{business logic} yang berat, serta interaksi \textit{multi-user concurrent} yang simultan \cite{maddodi2018generating}.

Kesenjangan ini menunjukkan bahwa belum ada bukti empiris yang memadai mengenai kelayakan \textit{deployment} ERP pada SBC: apakah perangkat ini mampu menangani beban ERP dari sisi performa, dan apakah \textit{Total Cost of Ownership}-nya benar-benar lebih hemat dibanding alternatif \textit{cloud} yang umum digunakan seperti VPS? Maka muncul pertanyaan penelitian: \textbf{"Bagaimana kelayakan \textit{deployment} ERP pada \textit{single-board computer} dibandingkan dengan VPS dari aspek performa dan biaya?"}

% --- Rumusan Masalah ---
\section{Rumusan Masalah}
\begin{enumerate}[1.]
    \item Bagaimana kelayakan \textit{deployment} ERP pada \textit{single-board computer} dibandingkan dengan VPS dari aspek performa dan biaya?
    \item Apa saja batasan \textit{deployment} ERP pada \textit{single-board computer}, khususnya dari aspek kapasitas \textit{user}, \textit{resource utilization}, dan kompleksitas modul?
\end{enumerate}

% --- Tujuan ---

\section{Tujuan}
\begin{enumerate}[1.]
    \item Melakukan studi kelayakan \textit{deployment} ERP pada \textit{single-board computer} dari aspek performa dan biaya.
    \item Melakukan perbandingan dan \textit{trade-off analysis} antara \textit{deployment} ERP pada \textit{single-board computer} dan VPS.
    \item Menemukan batasan \textit{deployment} ERP pada \textit{single-board computer} yang digunakan
 \end{enumerate}

% --- Batasan Masalah ---
\section{Batasan Masalah}
\begin{enumerate}[1.]
    \item Penelitian ini berfokus pada \textit{deployment} ERP di \textit{single-board computer} dengan perbandingan terhadap \textit{Virtual Private Server} (VPS) \textit{cloud}.
    \item Untuk evaluasi performa, penelitian menggunakan VPS dengan spesifikasi yang memenuhi atau melebihi \textit{minimum requirement} dari \textit{platform} ERP yang dipilih.
    \item Untuk analisis \textit{Total Cost of Ownership} (TCO), penelitian menggunakan \textit{commercial} VPS yang \textit{representative} untuk \textit{use case deployment} ERP, dengan model \textit{pricing} yang \textit{transparent} dan \textit{measurable} untuk proyeksi jangka panjang. Platform dengan biaya Rp 0 (\textit{free tier}) tidak dimasukkan dalam analisis kuantitatif TCO.
    \item ERP yang digunakan untuk perbandingan bersifat \textit{open-source} (contoh: Odoo Community atau ERPNext) agar tidak menambah biaya lisensi penelitian.
    \item Pengujian dibatasi pada modul esensial ERP (keuangan, inventori, CRM) dengan simulasi beban kerja yang representatif untuk organisasi skala kecil-menengah (10-50 \textit{concurrent users}).
    \item Penelitian ini mengkaji kelayakan teknis (performa sistem dalam menangani beban kerja) dan finansial (\textit{Total Cost of Ownership} dibandingkan dengan alternatif VPS), tidak mencakup aspek \textit{organizational readiness}, \textit{change management}, atau \textit{user training}.
    \item Batasan \textit{deployment} ERP yang diidentifikasi dalam penelitian ini bersifat spesifik terhadap perangkat SBC yang diuji, namun dapat menjadi indikasi untuk perangkat sejenis dengan spesifikasi \textit{hardware} yang \textit{comparable}.
\end{enumerate}

% --- Metodologi Pengerjaan TA ---
\section{Metode}
\begin{enumerate}
    \item \textbf{Studi Literatur} \\
      Mengumpulkan referensi mengenai ERP, \textit{deployment} pada SBC dan \textit{cloud}, serta metodologi analisis biaya kepemilikan (TCO).

    \item \textbf{Persiapan Infrastruktur} \\
      Menyiapkan SBC dan VPS dengan spesifikasi yang \textit{comparable}, kemudian melakukan instalasi \textit{platform} ERP \textit{open-source} (Odoo atau ERPNext) dengan konfigurasi identik.

    \item \textbf{Pengujian Performa} \\
      Melakukan \textit{load testing} dengan simulasi \textit{concurrent users} (10-50 \textit{users}) untuk mengukur \textit{response time}, \textit{throughput}, dan \textit{resource utilization} pada kedua infrastruktur.

    \item \textbf{Analisis \textit{Total Cost of Ownership} (TCO)} \\
      Menghitung dan membandingkan biaya kepemilikan untuk periode 1 tahun dan 3 tahun, mencakup biaya \textit{hardware}, operasional, dan konsumsi daya.

    \item \textbf{Analisis Hasil} \\
      Membandingkan hasil performa dan TCO, melakukan \textit{trade-off analysis}, dan mengidentifikasi batasan \textit{deployment} ERP pada SBC.
\end{enumerate}